% ===== 10 结论 =====
\section{结论}
\label{sec:conclusion}

本文面向混凝土骨料自然堆积场景的智能筛分任务，构建了一套从现场图像到颗粒实例、从像素几何到毫米尺度、再从二维视觉测量到质量级配输出的分层方法。与单纯将任务理解为“颗粒检测”不同，本文将赛题的关键难点归纳为三个连续问题：如何在粘连与接触条件下恢复实例，如何建立可靠物理尺度，以及如何把二维可见颗粒统计转换为与机械筛分相同的质量语义。围绕这三个问题，系统形成了“实例感知--物理测量--筛分语义映射--形貌/异常--自动报告”的完整处理链。

在视觉感知层，本文复用 ImageGrains 2.0 / Cellpose-SAM 的预训练实例分割能力，避免把有限现场时间消耗在重新标注和训练上；在物理测量层，通过显式尺度 $r$ 和实例轮廓提取长轴、短轴、面积及其他几何量，使每一个毫米级结果都能追溯到图像与尺度标定；在筛分语义层，定义
\begin{equation}
 d_i=\theta_1 b_i+\theta_2 d_{\mathrm{eq},i}+\theta_3,
 \qquad
 w_i=d_i^{\gamma},
\end{equation}
分别处理“二维投影尺寸不等于筛孔尺寸”和“颗粒数量不等于颗粒质量”两个问题，并通过累计加权分布得到 $D_{10}/D_{50}/D_{90}$ 与标准筛级质量占比。该结构的重点不是固定使用某一组经验参数，而是把视觉--筛分差异压缩为少量可解释参数，使其能够由同批次机械筛分真值校准并在独立批次上验证。

当前功能演示集包含三张真实自然堆积骨料图像，分别产生 1924、968 和 966 个实例，并能够完整生成逐颗粒几何表、预测级配、投影形貌、异常标签和可视化结果。三张样本均显示数量型与 $d^3$ 质量代理型分布存在明显差异，例如 \texttt{agg\_001} 的 $D_{50}$ 从数量口径约 6.0~mm 移动到质量代理口径约 23.5~mm。这一结果证明“数量统计不能直接当作质量筛分”是必须显式处理的建模问题，但并不单独证明 $\gamma=3$ 或当前 $d=b$ 已经达到机械筛分精度。仓库的 27 项合成竞赛层测试进一步验证了等效粒径、加权分位数、校准、形貌、异常和报告模块在软件定义上的一致性；已有处理记录表明算法计算阶段具有充足的 30~min 现场时间余量。

同时，本文没有回避当前最关键的证据缺口：三张演示样本尚未配对提交标准机械筛分逐筛质量，因此现阶段不能给出可信的 $D_p$ 预测误差、完整筛孔曲线误差或“筛分准确率”结论；投影形貌和疑似泥团规则同样缺少独立人工/物理标签，不能仅凭规则输出比例声称分类准确。基于这一认识，报告将功能验证与精度验证严格分离，并建立按物理批次划分开发、校准和独立测试数据的实验协议。正式竞赛版本的最高优先级，是采集同批图像--机械筛分配对数据，拟合并冻结 $\boldsymbol{\theta},\gamma$，随后在未参与校准的批次上报告 $D_{10}/D_{50}/D_{90}$、逐筛级质量占比和完整累计通过率误差。

从工程角度看，本文方案已经从“算法脚本”扩展为可部署流程：现场先完成稳定成像与尺度标定，再运行固定参数实例分割和筛分分析，通过 composite、轴长叠加、异常高亮和场景汇总完成质检，并保存原图、尺度、模型参数、逐颗粒 CSV、场景 JSON 与可视化结果形成追溯链。30~min 约束也被拆解为设备布置、标定、采集、计算、质检、导出和故障余量，而不是只比较单张 GPU 推理时间。

因此，本文当前已经完成的是一套\textbf{结构完整、可运行、可解释且预留物理校准闭环的视觉筛分系统}；尚需完成的是用独立标准筛分和人工标注把“可运行”进一步升级为“可量化证明准确”。这一划分也给出了明确的后续路线：首先补齐机械筛分与实例 ground truth，其次改善多点尺度/畸变校正和困难场景实例感知，最后扩展泥团材质分类与三维针片状测量。完成这些验证后，本文的分层框架无需改变，只需以真实校准参数和独立测试结果替换当前经验基线，即可形成面向比赛最终评分的完整技术报告。
